
\documentclass[11pt]{article}

\usepackage[margin=1in]{geometry}
\usepackage{amsmath, amssymb, amsthm}
\usepackage[T1]{fontenc}
\usepackage[utf8]{inputenc}
\usepackage{lmodern}
\usepackage{enumitem}
\usepackage{hyperref}
\usepackage{mathtools}

\title{Emergent Structure Without Conservation}
\author{}
\date{}

\begin{document}
\maketitle

\section*{Overview}

We construct a minimal dynamical system that:

\begin{itemize}[leftmargin=2em]
    \item has no conserved quantities,
    \item assumes no substances,
    \item only has:
    \begin{itemize}
        \item state,
        \item transformation,
        \item observables,
    \end{itemize}
    \item and still produces:
    \begin{itemize}
        \item structure,
        \item persistence,
        \item ``entities'' (via thresholds).
    \end{itemize}
\end{itemize}

\section*{1. Define the system (no ontology)}

Let the state be
\[
x_t \in \mathbb{R}^n
\]
and let the evolution be given by
\[
x_{t+1} = F(x_t).
\]

Choose a concrete example:
\[
x_{t+1} = \tanh(Ax_t) + \epsilon_t
\]
where:
\begin{itemize}[leftmargin=2em]
    \item \(A\) is a random matrix,
    \item \(\epsilon_t\) is noise.
\end{itemize}

Important properties:
\begin{itemize}[leftmargin=2em]
    \item nonlinear,
    \item stochastic,
    \item no conserved quantity,
    \item no symmetry assumed.
\end{itemize}

This is a pure transformation system.

\section*{2. No ``stuff,'' only observables}

Define observables:
\[
g_i(x) = \text{any function of } x.
\]

Examples:
\begin{align*}
g_1(x) &= x_1 \\
g_2(x) &= \|x\|^2 \\
g_3(x) &= x_1 x_2
\end{align*}

These are not substances. They are just ways of reading the system.

\section*{3. Generate ``flows''}

Each observable produces a time series:
\[
g_i(x_t).
\]

This is the sense in which we use the word ``flow'' here:
\begin{itemize}[leftmargin=2em]
    \item not something moving as a substance,
    \item but a trajectory of a measurement.
\end{itemize}

\section*{4. Extract structure (no conservation needed)}

Now we detect structure using the following criteria.

\subsection*{A. Persistence}

Find regions where the system stays for long times.

\subsection*{B. Predictability}

Find observables \(g_i\) such that
\[
g_i(x_{t+1}) \approx f(g_i(x_t)).
\]

\subsection*{C. Slow modes}

Find directions where change is slow.

These correspond to:
\begin{itemize}[leftmargin=2em]
    \item metastability,
    \item Koopman modes,
    \item predictive coarse-graining.
\end{itemize}

\section*{5. Define ``entities''}

Now define:
\[
E = \{x \mid g_i(x) > c\}.
\]

But choose \(g_i\) and \(c\) by:
\begin{itemize}[leftmargin=2em]
    \item persistence,
    \item predictability.
\end{itemize}

Result:
\[
\text{entities} = \text{stable regions of observable space}.
\]

\section*{6. Iteration (self-referential part)}

Now:
\begin{enumerate}[leftmargin=2em]
    \item choose observables \(g_i\),
    \item evaluate:
    \begin{itemize}
        \item predictability,
        \item persistence,
    \end{itemize}
    \item update observables.
\end{enumerate}

Formally:
\[
g^{(k+1)} = \mathcal{A}\bigl(g^{(k)}\bigr).
\]

This is selection without conservation.

\section*{7. What emerges}

Even though we assumed:
\begin{itemize}[leftmargin=2em]
    \item no conserved quantities,
    \item no substances,
\end{itemize}
we still get:
\begin{itemize}[leftmargin=2em]
    \item clusters,
    \item regimes,
    \item slow dynamics,
    \item stable regions.
\end{itemize}

These look like:
\begin{itemize}[leftmargin=2em]
    \item ``entities,''
    \item ``flows.''
\end{itemize}

\section*{8. Important insight}

Nothing was conserved.

Nothing was fundamental.

Yet structure emerged anyway.

\section*{9. Interpretation}

This shows:
\begin{itemize}[leftmargin=2em]
    \item conservation is not required,
    \item substance is not required.
\end{itemize}

What matters is:
\[
\text{structure in dynamics} + \text{selection over observables}.
\]

\section*{10. Clean ontology}

We can now state it rigorously:
\begin{itemize}[leftmargin=2em]
    \item system \(=\) transformation \(F\),
    \item observables \(=\) functions \(g(x)\),
    \item flows \(=\) trajectories \(g(x_t)\),
    \item entities \(=\) thresholded regions of \(g\),
    \item significance \(=\) persistence \(+\) predictability.
\end{itemize}

\section*{11. What this proves}

You can:
\begin{itemize}[leftmargin=2em]
    \item build a model,
    \item extract structure,
    \item define entities,
\end{itemize}
without ever invoking:
\begin{itemize}[leftmargin=2em]
    \item conserved quantities,
    \item substances,
    \item fixed ontology.
\end{itemize}

\section*{Final statement}

This is the clean version of the idea:

\begin{quote}
A world can emerge from pure dynamics and selection over observables, without any conserved ``stuff'' or predefined entities.
\end{quote}

\end{document}
